\clearpage
\fancyhead[L]{\small MÜLLER. CALCULATION OF SOLAR ECLIPSES.}
\hypertarget{native-596-lower}{}
\begin{center}\small Göttingische gelehrte Anzeigen. Issue 20. Page 200. 3 February 1816.\\[3pt]\rule{.72\linewidth}{.3pt}\end{center}

\emph{Detailed Guide to the Trigonometric Calculation of Solar Eclipses Visible at a Given Place on the Earth’s Surface}, by two different, very accurate methods, illustrated by determining the phenomena of the great solar eclipse of 19 November 1816 for the Nuremberg meridian, by \textsc{Johann Wolfgang Müller}, Professor of Mathematics at the Royal Gymnasium in Nuremberg. Sulzbach. At the art and book dealership of Commercial Councillor J. E. \textsc{Seidel}. 128 pages in octavo, with a copperplate.

The author has taken great care to calculate in advance the principal phases of this year’s remarkable solar eclipse for the Nuremberg horizon, where its magnitude will be 10 digits and 38 minutes. He used the trigonometric method of the nonagesimal, first following \textsc{Tob. Mayer}’s arrangement exactly and then following a slightly modified course of calculation, with no real change of substance. He has had all the operations printed in this little work, down to the smallest detail. It can be recommended as a practical, clear guide for beginners and inexperienced amateurs who would like to carry out a similar calculation for the place where they live.
